% Sorgente LaTeX canonico, importato e revisionato dal precedente sorgente Markdown.
\chapter{Funzioni Gamma e Beta}
\label{chap:19}

\begin{obiettivocapitolo}{ricondurre integrali speciali alle funzioni Gamma e Beta}
\estensione. Le formule sono utili per integrali, probabilità e normalizzazioni.
\end{obiettivocapitolo}
\begin{sceltametodo}{riconoscimento}
Una potenza di $t$ moltiplicata per $e^{-t}$ su $(0,\infty)$ suggerisce Gamma. Un prodotto $t^{p-1}(1-t)^{q-1}$ su $(0,1)$ suggerisce Beta. Per esponenziali $e^{-a x^r}$ usa una sostituzione che produca $t=a x^r$.
\end{sceltametodo}
\begin{proceduraesame}{riduzione}
\begin{enumerate}
\item Controlla le condizioni sui parametri che garantiscono convergenza.
\item Esegui una sostituzione per ottenere esattamente una forma standard.
\item Applica ricorrenza, relazione Beta--Gamma o valori notevoli.
\item Semplifica fattoriali e potenze mantenendo le costanti di scala.
\end{enumerate}
\end{proceduraesame}
\begin{controllofinale}{integrali speciali}
Verifica positività quando l'integrando è positivo, dimensione dei parametri e casi interi/seminteri noti.
\end{controllofinale}

Argomento di approfondimento, utile soprattutto nelle applicazioni scientifiche e probabilistiche. Nelle definizioni, \(t\) è una variabile muta di integrazione. La funzione Gamma estende il fattoriale dagli interi positivi ai reali positivi.

\section{Definizioni e identità fondamentali}\label{definizioni-e-identituxe0-fondamentali}

Per \(x>0\):

\[
\displaystyle
\Gamma(x)
=
\int_0^{+\infty}t^{x-1}e^{-t}\,dt,
\]

\[
\displaystyle
\Gamma(x+1)=x\Gamma(x),
\qquad
\Gamma(n)=(n-1)!\quad(n\in\mathbb N_+),
\qquad
\Gamma(n+1)=n!\quad(n\in\mathbb N_0),
\qquad
\Gamma\left(\frac12\right)=\sqrt\pi.
\]

Per \(x,y>0\):

\[
\displaystyle
B(x,y)
=
\int_0^1 t^{x-1}(1-t)^{y-1}\,dt,
\]

\[
\displaystyle
B(x,y)=B(y,x)
=
\frac{\Gamma(x)\Gamma(y)}{\Gamma(x+y)},
\]

\[
\displaystyle
B(x+1,y)=\frac{x}{x+y}B(x,y),
\qquad
B(x,y+1)=\frac{y}{x+y}B(x,y).
\]

Teoria: \href{https://ingegnerismo.it/matematica/funzione-gamma/}{funzione Gamma} e \href{https://ingegnerismo.it/matematica/funzione-beta/}{funzione Beta}.
